\documentclass[11pt]{article}

\usepackage[margin=1in]{geometry}
\usepackage{amsmath,amssymb,amsthm}
\usepackage{booktabs}
\usepackage{hyperref}

\newtheorem{theorem}{Theorem}
\newtheorem{proposition}[theorem]{Proposition}
\newtheorem{corollary}[theorem]{Corollary}
\newtheorem{remark}[theorem]{Remark}

\title{Remedies for the Force-Magnitude Dip in an\\
       Inverse-Elliptical Potential Field}
\author{}
\date{}

\begin{document}
\maketitle

\begin{abstract}
A companion to \emph{Why an Inverse-Elliptical Potential Produces a
Force-Magnitude Dip Along its Major Axis}, which established that
$|\nabla U|$ has a strict local minimum on the major axis whenever the
elongation $q$ exceeds $\sqrt{3}$. Here we identify precisely which factor of
the gradient carries the defect, and evaluate four candidate remedies against
two requirements: the dip must be removed, and the force \emph{direction}
must be preserved so that the tangential component driving lateral evasion
survives. We show that normalising the gradient by $|\nabla\rho|$ is the only
candidate meeting both, prove it removes the dip for every $q$, and quantify
the price paid, namely loss of conservativeness. A decision rule in terms of
$q$ concludes.
\end{abstract}

\section{What is and is not the defect}
\label{sec:defect}

It is worth first excluding a plausible but incorrect diagnosis. On the major
axis one has $F_{x}\equiv0$ whenever the obstacle is laterally symmetric, since
$U(x,y)=U(-x,y)$ implies $\partial_{x}U(0,y)=0$ identically. This is a
symmetry statement, independent of $q$, and it holds equally for the
\emph{circular} field $q=1$, which exhibits no dip at all. Vanishing lateral
force on the axis is therefore not the phenomenon under study and cannot be
what distinguishes the elliptical case.

The genuine distinction is the sign of the second derivative, i.e.\ the
character of the critical point. Along a horizontal scan at height $s$:
\begin{itemize}
  \item for $q<\sqrt{3}$ (including the circle), $x=0$ is a \emph{maximum} of
        $|F|$, which then decays monotonically outward;
  \item for $q>\sqrt{3}$, $x=0$ is a \emph{minimum}, and $|F|$ first rises to
        an off-axis ridge before decaying.
\end{itemize}
The force is weakest exactly where the obstacle lies most directly ahead. The
practical consequence appears whenever $|F|$ is thresholded as a scalar risk
measure: at $q=12$, $s=9.5$ the on-axis value is $4.80$ while the flank ridge
reaches $12.12$, so a trigger set at $8$ fires on the flanks but not straight
ahead.

\begin{remark}[The dip is relative, not absolute]
\label{rem:relative}
Elongation does not weaken the on-axis force. With $t_{x}=1$ and
$q=t_{y}/t_{x}$, the elliptical radius on the axis is $\rho(0,s)=s/t_{y}$, so
$U(0,s)=\lambda t_{y}/(s-l)$ and
\begin{equation}
  \bigl|F_{y}(0,s)\bigr|=\frac{\lambda\,t_{y}}{(s-l)^{2}},
  \label{eq:onaxis}
\end{equation}
which \emph{increases} linearly with $t_{y}$; numerically
$1.200,\,2.000,\,4.800$ for $q=3,5,12$ at $s=9.5$, $\lambda=10$, $l=4.5$.
The ridge merely grows faster, as $O(t_{y}^{2})$. The dip is thus a purely
relative feature, and any remedy should be judged on restoring monotonicity,
not on raising the on-axis magnitude.
\end{remark}

\section{Localising the defect}

Let
\begin{equation}
  \rho=\sqrt{x^{2}+\frac{y^{2}}{q^{2}}},\qquad U=U(\rho),
\end{equation}
so that $U$ depends on position only through $\rho$. By the chain rule
$\nabla U=U'(\rho)\nabla\rho$, hence
\begin{equation}
  \bigl|\nabla U\bigr|
  =\underbrace{\bigl|U'(\rho)\bigr|}_{\text{function of }\rho\text{ alone}}
  \cdot
  \underbrace{\bigl|\nabla\rho\bigr|}_{\text{varies around the ellipse}} .
  \label{eq:factorisation}
\end{equation}

\begin{proposition}[The first factor is harmless]
\label{prop:harmless}
At fixed height $s>0$, $\rho(x,s)$ is strictly increasing in $|x|$. Hence if
$U'$ is of one sign and $|U'|$ is monotone in $\rho$, the factor $|U'(\rho)|$
is monotone in $|x|$ with its extremum at $x=0$, and can never produce an
interior ridge.
\end{proposition}

\begin{proof}
$\partial_{x}\rho=x/\rho>0$ for $x>0$, so $\rho$ strictly increases with
$|x|$; composing a monotone function of $\rho$ with a monotone function of
$|x|$ preserves monotonicity.
\end{proof}

The whole defect therefore resides in $|\nabla\rho|$. Explicitly,
\begin{equation}
  \bigl|\nabla\rho\bigr|=\frac{\sqrt{x^{2}+y^{2}/q^{4}}}{\rho},
  \qquad
  \bigl|\nabla\rho\bigr|
  =\begin{cases}
     1/q & \text{on the major axis } (x=0),\\[2pt]
     1   & \text{on the minor axis } (y=0).
   \end{cases}
  \label{eq:gradrho-range}
\end{equation}
This factor sweeps the range $[1/q,\,1]$ as one moves around an ellipse, and
it is precisely this $q$-fold swing that overturns the decay of $|U'|$ once
$q^{2}>3$. Equation~\eqref{eq:gradrho-range} immediately suggests the two
lines of attack pursued below: \emph{shrink the range} of the offending
factor (reduce $q$), or \emph{divide it out}.

\section{Candidate remedies}

We evaluate each candidate against two requirements:

\begin{description}
  \item[(R1)] the contrast $C=|F|_{\max}/|F|_{x=0}$ on a horizontal scan
        equals $1$, i.e.\ the dip is removed;
  \item[(R2)] the force \emph{direction} is preserved, so the tangential
        component that drives lateral evasion is not destroyed.
\end{description}

Requirement (R2) is not cosmetic. A remedy that leaves the force purely
radial reproduces the classical artificial-potential-field deadlock, in which
a vehicle directly behind an obstacle receives braking but no lateral
guidance.

\subsection{Remedy A: reduce the elongation \texorpdfstring{$q$}{q}}

By Theorem~1 of the companion note, the dip vanishes outright for
$q\le\sqrt3$. In the implemented field the effective threshold is more
forgiving, because the vehicle-dimension prefactor $(1-l/s)^{-1}$ contributes
an additional $x$-dependent term to $F_{y}$ that partially fills the minimum.
Measured on the implemented form:

\begin{center}
\begin{tabular}{rccc}
\toprule
$q=t_{y}/t_{x}$ & $C$, pure ellipse & $C$ at $s=7.5$ & $C$ at $s=9.5$\\
\midrule
$12$        & $4.651$ & $1.982$ & $2.524$\\
$5$         & $2.005$ & $1.091$ & $1.241$\\
$3$         & $1.299$ & $1.000$ & $1.007$\\
$2$         & $1.026$ & $1.000$ & $1.000$\\
$\sqrt3$    & $1.000$ & $1.000$ & $1.000$\\
$1$         & $1.000$ & $1.000$ & $1.000$\\
\bottomrule
\end{tabular}
\end{center}

\noindent
Thus $q\le3$ already satisfies (R1) to three decimal places in practice.
Requirement (R2) is met trivially, the field being unmodified. Remedy~A
preserves conservativeness and retains $U$ as a candidate Lyapunov function;
its only cost is that $q$ is a modelling parameter, so lowering it narrows
the longitudinal reach of the danger zone.

\subsection{Remedy B: steepen the radial power}

Generalise to $U=\rho^{-n}$. Then $\nabla U=-n\rho^{-n-2}(x,\,y/q^{2})$ and,
with $u=x^{2}$ and $a=s/q$,
\begin{equation}
  g(u)=|\nabla U|^{2}\big|_{y=s}
      =n^{2}\,\frac{u+a^{2}/q^{2}}{(u+a^{2})^{\,n+2}} .
\end{equation}
Differentiating and cancelling $(u+a^{2})^{n+1}$ as before,
\begin{equation}
  \operatorname{sign}g'(u)
  =\operatorname{sign}\Bigl[\,a^{2}\Bigl(1-\frac{n+2}{q^{2}}\Bigr)-(n+1)u\,\Bigr],
\end{equation}
so the dip criterion generalises to
\begin{equation}
  \boxed{\;q>\sqrt{n+2}\;}
\end{equation}
recovering $\sqrt3$ at $n=1$. Raising $n$ does raise the threshold, but far
too slowly to be useful, and the response is not even monotone in $n$:

\begin{center}
\begin{tabular}{rccc}
\toprule
$n$ & $\sqrt{n+2}$ & $C$ at $q=5$ & $C$ at $q=12$\\
\midrule
$1$  & $1.73$ & $1.241$ & $2.524$\\
$2$  & $2.00$ & $1.302$ & $2.760$\\
$3$  & $2.24$ & $1.280$ & $2.712$\\
$6$  & $2.83$ & $1.170$ & $2.390$\\
$10$ & $3.46$ & $1.074$ & $2.062$\\
$23$ & $5.00$ & $1.000$ & $1.549$\\
\bottomrule
\end{tabular}
\end{center}

\noindent
Suppressing the dip at $q=5$ requires $n\approx23$, and $n=2$ is worse than
$n=1$. Remedy~B is rejected.

\subsection{Remedy C: the metric-consistent force}

Take $F=-G^{-1}\nabla U$ with $G=\operatorname{diag}(1,1/q^{2})$, the metric
induced by $\rho$. For $U=1/\rho$ this gives $F=(x,y)\rho^{-3}$, so
$|F|=r\rho^{-3}$ with $r=\sqrt{x^{2}+y^{2}}$, and
\begin{equation}
  g(u)=\frac{u+s^{2}}{(u+a^{2})^{3}},
  \qquad
  \operatorname{sign}g'(u)=\operatorname{sign}\bigl[a^{2}-3s^{2}-2u\bigr]<0
\end{equation}
for all $u\ge0$ whenever $a<s\sqrt3$, i.e.\ $q>1/\sqrt3$. Requirement (R1) is
met for every practical $q$.

Requirement (R2), however, fails. The force becomes exactly radial, and the
measured force angle at $x=0.5$, $s=9.5$ collapses to $1.6^{\circ}$,
$1.5^{\circ}$, $1.3^{\circ}$, $0.7^{\circ}$ for $q=3,5,12,30$ --- against
$13.8^{\circ}$, $33.8^{\circ}$, $73.4^{\circ}$, $85.0^{\circ}$ for the
unmodified field. Essentially all lateral guidance is destroyed. Remedy~C is
rejected.

\subsection{Remedy D: normalise by \texorpdfstring{$|\nabla\rho|$}{|grad rho|}}

Equation~\eqref{eq:factorisation} localised the defect in $|\nabla\rho|$;
Remedy~D simply divides it out. Define
\begin{equation}
  F \;=\; -\,\frac{\nabla U}{\bigl|\nabla\rho\bigr|}
    \;=\; -\,U'(\rho)\,\hat{n},
  \qquad
  \hat{n}=\frac{\nabla\rho}{\bigl|\nabla\rho\bigr|} .
  \label{eq:remedyD}
\end{equation}

\begin{theorem}[Remedy D removes the dip at every elongation]
\label{thm:D}
For $F$ as in \eqref{eq:remedyD} with $U=1/\rho$, and for every $q>0$ and
every height $s>0$, the magnitude $|F|$ is strictly decreasing in $|x|$ and
attains its maximum at $x=0$. In particular $C=1$, and no off-axis ridge
exists.
\end{theorem}

\begin{proof}
Since $\hat{n}$ is a unit vector, $|F|=|U'(\rho)|=\rho^{-2}$, a function of
$\rho$ alone and strictly decreasing in $\rho$. At fixed $s$,
$\rho(x,s)=\sqrt{x^{2}+(s/q)^{2}}$ is strictly increasing in $|x|$ by
Proposition~\ref{prop:harmless}. The composition of a strictly decreasing
with a strictly increasing function is strictly decreasing, so $|F|$ attains
its maximum at $x=0$ and has no interior critical point. No hypothesis on $q$
was used.
\end{proof}

\begin{corollary}[Direction is untouched]
Dividing by the strictly positive scalar $|\nabla\rho|$ rescales the vector
without rotating it, so $F$ in \eqref{eq:remedyD} is everywhere parallel to
$-\nabla U$. Requirement (R2) holds exactly.
\end{corollary}

\noindent
Both claims are confirmed numerically at $s=9.5$. The angle column is
identical to that of the unmodified field, and the contrast is exactly unity:

\begin{center}
\begin{tabular}{lrccccc}
\toprule
mode & $q$ & $|F|(0)$ & peak & $x^{*}$ & $C$ & angle at $x=0.5$\\
\midrule
plain  & $3$  & $1.200$  & $1.208$  & $0.828$ & $1.007$ & $13.8^{\circ}$\\
plain  & $5$  & $2.000$  & $2.482$  & $1.070$ & $1.241$ & $33.8^{\circ}$\\
plain  & $12$ & $4.800$  & $12.117$ & $0.540$ & $2.524$ & $73.4^{\circ}$\\
plain  & $30$ & $12.000$ & $73.376$ & $0.223$ & $6.115$ & $85.0^{\circ}$\\
\midrule
norm   & $3$  & $3.600$   & $3.600$   & $0$ & $1.000$ & $13.8^{\circ}$\\
norm   & $5$  & $10.000$  & $10.000$  & $0$ & $1.000$ & $33.8^{\circ}$\\
norm   & $12$ & $57.600$  & $57.600$  & $0$ & $1.000$ & $73.4^{\circ}$\\
norm   & $30$ & $360.000$ & $360.000$ & $0$ & $1.000$ & $85.0^{\circ}$\\
\midrule
metric & $3$  & $10.800$    & $10.800$    & $0$ & $1.000$ & $1.6^{\circ}$\\
metric & $5$  & $50.000$    & $50.000$    & $0$ & $1.000$ & $1.5^{\circ}$\\
metric & $12$ & $691.200$   & $691.200$   & $0$ & $1.000$ & $1.3^{\circ}$\\
metric & $30$ & $10800.000$ & $10800.000$ & $0$ & $1.000$ & $0.7^{\circ}$\\
\bottomrule
\end{tabular}
\end{center}

\section{The price of Remedy D}

\paragraph{Scale.}
Since $|\nabla\rho|\to1/q$ on the axis, the normalisation inflates magnitudes
by up to a factor $q$ there ($4.800\to57.600$ at $q=12$). If absolute force
values are calibrated, rescale $\lambda$ by $O(1/q)$.

\paragraph{Loss of conservativeness.}
Write $F=-\,\Phi\,\nabla\rho$ with $\Phi=U'(\rho)/|\nabla\rho|$. A field of
this form is a gradient if and only if $\Phi$ is a function of $\rho$ alone.
But by \eqref{eq:gradrho-range}, $|\nabla\rho|$ varies from $1/q$ to $1$
\emph{along a single ellipse} $\rho=\text{const}$, so $\Phi$ is not such a
function and $F$ is not conservative for any $q>1$. Consequently no scalar
potential has $F$ as its descent direction, and $U$ is no longer available as
a Lyapunov function.

The magnitude of the violation, measured as
$(\nabla\times F)/|F|$ at $(x,y-y_{o})=(1.0,\,7.5)$, is

\begin{center}
\begin{tabular}{rcc}
\toprule
$q$ & $(\nabla\times F)/|F|$ (m$^{-1}$) & $|F|$\\
\midrule
$1$  & $\phantom{-}0.0000$ & $1.095$\\
$3$  & $-0.4384$ & $6.603$\\
$5$  & $-0.4789$ & $10.163$\\
$12$ & $-0.1225$ & $18.907$\\
$30$ & $-0.0203$ & $24.095$\\
\bottomrule
\end{tabular}
\end{center}

\noindent
The vanishing entry at $q=1$ is a consistency check: for a circle
$|\nabla\rho|\equiv1$, so Remedy~D reduces to the identity. Note the
violation is \emph{non-monotone}, peaking near $q\approx5$ and becoming mild
at large $q$ --- fortunately, the regime in which Remedy~D is actually needed.

\section{Decision rule}

\begin{center}
\begin{tabular}{lcccl}
\toprule
Remedy & (R1) & (R2) & conservative & verdict\\
\midrule
A. reduce $q$ to $\le3$        & yes & yes & yes & preferred when $q$ may be lowered\\
B. steepen power $n$           & no  & yes & yes & rejected: needs $n\approx23$\\
C. metric gradient             & yes & no  & yes & rejected: radial, causes deadlock\\
D. normalise by $|\nabla\rho|$ & yes & yes & no  & preferred when large $q$ is required\\
\bottomrule
\end{tabular}
\end{center}

\noindent
Concretely:
\begin{itemize}
  \item \textbf{If $q\lesssim5$:} apply Remedy~A. At $q=5$ the contrast is
        already only $1.241$, and at $q=3$ it is $1.007$ --- the dip is gone
        for practical purposes, at no structural cost.
  \item \textbf{If $q\gtrsim12$ is required by the model:} apply Remedy~D. It
        is the only candidate satisfying both (R1) and (R2), and its
        circulation penalty is mildest precisely in that regime
        ($0.12\,$m$^{-1}$ at $q=12$ versus $0.48\,$m$^{-1}$ at $q=5$).
\end{itemize}

\noindent
Finally, a scale check that may render the question moot. The ridge sits at
$x^{*}\approx s/(q\sqrt2)$, so at $q=12$, $s=9.5$ one finds $x^{*}=0.554$\,m,
which is smaller than the half-width $0.875$\,m of a $1.75$\,m vehicle. The
entire low-force channel is then narrower than the host vehicle itself: a
vehicle centred on the axis already has its corners in the high-force ridge.
At such elongations the dip is largely an artefact of evaluating a point
field, and no remedy may be warranted at all.

\end{document}
